\documentclass[12pt,a4paper]{article}
\usepackage[margin=1in]{geometry}
\usepackage{listings}
\usepackage{xcolor}
\usepackage{hyperref}
\usepackage{fancyhdr}
\usepackage{array}
\usepackage{booktabs}

% Code listing setup
\lstdefinelanguage{VHDL}{
  morekeywords={library, use, entity, port, in, out, is, end, architecture, signal, 
                process, if, then, elsif, else, case, when, others, rising_edge},
  morecomment=[l]{--},
  basicstyle=\ttfamily\small,
  keywordstyle=\color{blue}\bfseries,
  commentstyle=\color{gray}\itshape,
  breaklines=true,
  showstringspaces=false,
  tabsize=2,
  numbers=left,
  numberstyle=\tiny\color{gray},
  backgroundcolor=\color{lightgray!10},
  frame=single,
}

\pagestyle{fancy}
\fancyhf{}
\rhead{VHDL Conditional Statements}
\lhead{Digital System Design}
\cfoot{\thepage}

\title{\textbf{VHDL Conditional Statements}\\
        \normalsize Case, If-Else, and When-Else}
\author{Digital System Design Guide}
\date{\today}

\begin{document}

\maketitle
\tableofcontents
\newpage

\section{Introduction}

VHDL provides three primary mechanisms for implementing conditional logic:
\begin{enumerate}
    \item \textbf{Case Statements} - For multiple mutually exclusive conditions
    \item \textbf{If-Else Statements} - For hierarchical conditional logic
    \item \textbf{When-Else Statements} - For concurrent signal assignments
\end{enumerate}

\subsection{When to Use Each Statement}

\begin{table}[h]
\centering
\begin{tabular}{|l|p{4cm}|p{4cm}|}
\hline
\textbf{Statement} & \textbf{Best Used For} & \textbf{Context} \\
\hline
\textbf{Case} & Multiple exclusive conditions & Sequential (process) \\
\hline
\textbf{If-Else} & Hierarchical/nested conditions & Sequential (process) \\
\hline
\textbf{When-Else} & Concurrent signal assignment & Concurrent \\
\hline
\end{tabular}
\end{table}

\newpage
\section{IF-ELSE Conditional Statement}

\subsection{Basic Structure}

\begin{lstlisting}[language=VHDL]
if (condition1) then
    -- statements executed if condition1 is TRUE
    
elsif (condition2) then
    -- statements executed if condition2 is TRUE
    
else
    -- statements executed if all conditions are FALSE
end if;
\end{lstlisting}

\subsection{Key Characteristics}
\begin{itemize}
    \item Used within \textbf{sequential processes}
    \item Supports \textbf{hierarchical/nested conditions}
    \item Priority is determined by \textbf{order of conditions}
    \item \textbf{Only one branch} executes
\end{itemize}

\subsection{Example 1: D-Latch}

\begin{lstlisting}[language=VHDL]
library IEEE;
use IEEE.STD_LOGIC_1164.all;

entity D_Latch is
    port (
        Enable  : in  STD_LOGIC;
        Data_In : in  STD_LOGIC;
        Q       : out STD_LOGIC;
        Q_bar   : out STD_LOGIC
    );
end D_Latch;

architecture Behavioral of D_Latch is
    signal Q_internal : STD_LOGIC := '0';
begin
    
    process(Enable, Data_In)
    begin
        if (Enable = '1') then
            -- Output follows input when enabled
            Q_internal <= Data_In;
        -- else: Q_internal holds previous value
        end if;
    end process;
    
    Q     <= Q_internal;
    Q_bar <= not Q_internal;
    
end Behavioral;
\end{lstlisting}

\subsubsection{Truth Table}

\begin{table}[h]
\centering
\begin{tabular}{|c|c|c|c|}
\hline
\textbf{Enable} & \textbf{Data In} & \textbf{Q} & \textbf{Behavior} \\
\hline
1 & 0 & 0 & Output follows input \\
\hline
1 & 1 & 1 & Output follows input \\
\hline
0 & X & Previous & Output holds \\
\hline
\end{tabular}
\end{table}

\newpage
\subsection{Example 2: D-Flip-Flop}

\begin{lstlisting}[language=VHDL]
library IEEE;
use IEEE.STD_LOGIC_1164.all;

entity D_FlipFlop is
    port (
        CLK     : in  STD_LOGIC;
        Reset   : in  STD_LOGIC;
        Data_In : in  STD_LOGIC;
        Q       : out STD_LOGIC;
        Q_bar   : out STD_LOGIC
    );
end D_FlipFlop;

architecture Behavioral of D_FlipFlop is
    signal Q_internal : STD_LOGIC := '0';
begin
    
    process(CLK, Reset)
    begin
        -- Asynchronous reset (highest priority)
        if (Reset = '1') then
            Q_internal <= '0';
            
        -- Synchronous data capture (rising edge)
        elsif rising_edge(CLK) then
            Q_internal <= Data_In;
        end if;
    end process;
    
    Q     <= Q_internal;
    Q_bar <= not Q_internal;
    
end Behavioral;
\end{lstlisting}

\subsubsection{Truth Table}

\begin{table}[h]
\centering
\begin{tabular}{|c|c|c|c|}
\hline
\textbf{CLK} & \textbf{Reset} & \textbf{Data In} & \textbf{Q} \\
\hline
$\uparrow$ & 0 & 0 & 0 \\
\hline
$\uparrow$ & 0 & 1 & 1 \\
\hline
X & 1 & X & 0 \\
\hline
\end{tabular}
\end{table}

\newpage
\section{CASE Statement}

\subsection{Basic Structure}

\begin{lstlisting}[language=VHDL]
case (select_signal) is
    when value1 =>
        -- statements if select_signal = value1
    
    when value2 =>
        -- statements if select_signal = value2
    
    when others =>
        -- default case
end case;
\end{lstlisting}

\subsection{Key Characteristics}
\begin{itemize}
    \item Used within \textbf{sequential processes}
    \item Evaluates a \textbf{single expression}
    \item \textbf{Mutually exclusive} branches
    \item Synthesizes to \textbf{MUX or decoder} logic
\end{itemize}

\subsection{Example 1: 4-to-1 Multiplexer}

\begin{lstlisting}[language=VHDL]
library IEEE;
use IEEE.STD_LOGIC_1164.all;

entity Mux4to1 is
    port (
        I0 : in  STD_LOGIC;
        I1 : in  STD_LOGIC;
        I2 : in  STD_LOGIC;
        I3 : in  STD_LOGIC;
        S  : in  STD_LOGIC_VECTOR(1 downto 0);
        Y  : out STD_LOGIC
    );
end Mux4to1;

architecture Behavioral of Mux4to1 is
begin
    
    process(I0, I1, I2, I3, S)
    begin
        case S is
            when "00" => Y <= I0;
            when "01" => Y <= I1;
            when "10" => Y <= I2;
            when "11" => Y <= I3;
            when others => Y <= '0';
        end case;
    end process;
    
end Behavioral;
\end{lstlisting}

\subsubsection{Truth Table}

\begin{table}[h]
\centering
\begin{tabular}{|c|c|}
\hline
\textbf{S[1:0]} & \textbf{Output Y} \\
\hline
00 & I0 \\
\hline
01 & I1 \\
\hline
10 & I2 \\
\hline
11 & I3 \\
\hline
\end{tabular}
\end{table}

\newpage
\subsection{Example 2: 2-to-4 Decoder}

\begin{lstlisting}[language=VHDL]
library IEEE;
use IEEE.STD_LOGIC_1164.all;

entity Decoder2to4 is
    port (
        I       : in  STD_LOGIC_VECTOR(1 downto 0);
        Enable  : in  STD_LOGIC;
        O       : out STD_LOGIC_VECTOR(3 downto 0)
    );
end Decoder2to4;

architecture Behavioral of Decoder2to4 is
begin
    
    process(I, Enable)
    begin
        O <= "0000";  -- Initialize
        
        if (Enable = '1') then
            case I is
                when "00" => O <= "0001";
                when "01" => O <= "0010";
                when "10" => O <= "0100";
                when "11" => O <= "1000";
                when others => O <= "0000";
            end case;
        end if;
    end process;
    
end Behavioral;
\end{lstlisting}

\subsubsection{Truth Table}

\begin{table}[h]
\centering
\begin{tabular}{|c|c|c|c|c|c|}
\hline
\textbf{I[1]} & \textbf{I[0]} & \textbf{O3} & \textbf{O2} & \textbf{O1} & \textbf{O0} \\
\hline
0 & 0 & 0 & 0 & 0 & 1 \\
\hline
0 & 1 & 0 & 0 & 1 & 0 \\
\hline
1 & 0 & 0 & 1 & 0 & 0 \\
\hline
1 & 1 & 1 & 0 & 0 & 0 \\
\hline
\end{tabular}
\end{table}

\newpage
\section{WHEN-ELSE Concurrent Assignment}

\subsection{Basic Structure}

\begin{lstlisting}[language=VHDL]
output_signal <= value1 when condition1 else
                 value2 when condition2 else
                 default_value;
\end{lstlisting}

\subsection{Key Characteristics}
\begin{itemize}
    \item Used for \textbf{concurrent signal assignment}
    \item Evaluates \textbf{multiple conditions}
    \item Priority determined by \textbf{order}
    \item \textbf{No process} required
\end{itemize}

\subsection{Example 1: Multiplexer with WHEN-ELSE}

\begin{lstlisting}[language=VHDL]
library IEEE;
use IEEE.STD_LOGIC_1164.all;

entity Mux4to1_WhenElse is
    port (
        I0 : in  STD_LOGIC;
        I1 : in  STD_LOGIC;
        I2 : in  STD_LOGIC;
        I3 : in  STD_LOGIC;
        S  : in  STD_LOGIC_VECTOR(1 downto 0);
        Y  : out STD_LOGIC
    );
end Mux4to1_WhenElse;

architecture Behavioral of Mux4to1_WhenElse is
begin
    
    Y <= I0 when (S = "00") else
         I1 when (S = "01") else
         I2 when (S = "10") else
         I3;
    
end Behavioral;
\end{lstlisting}

\newpage
\subsection{Example 2: Decoder with WHEN-ELSE}

\begin{lstlisting}[language=VHDL]
library IEEE;
use IEEE.STD_LOGIC_1164.all;

entity Decoder2to4_WhenElse is
    port (
        I       : in  STD_LOGIC_VECTOR(1 downto 0);
        Enable  : in  STD_LOGIC;
        O       : out STD_LOGIC_VECTOR(3 downto 0)
    );
end Decoder2to4_WhenElse;

architecture Behavioral of Decoder2to4_WhenElse is
begin
    
    O <= "0001" when (Enable = '1' and I = "00") else
         "0010" when (Enable = '1' and I = "01") else
         "0100" when (Enable = '1' and I = "10") else
         "1000" when (Enable = '1' and I = "11") else
         "0000";
    
end Behavioral;
\end{lstlisting}

\newpage
\section{4-Bit Serial-to-Parallel Shift Register}

\subsection{Overview}

A serial-to-parallel shift register converts serial input into parallel output by shifting data on each clock pulse.

\begin{lstlisting}[language=VHDL]
library IEEE;
use IEEE.STD_LOGIC_1164.all;

entity Serial_to_Parallel_SR is
    port (
        CLK       : in  STD_LOGIC;
        RESET     : in  STD_LOGIC;
        Shift     : in  STD_LOGIC;
        Serial_In : in  STD_LOGIC;
        S         : out STD_LOGIC_VECTOR(3 downto 0)
    );
end Serial_to_Parallel_SR;

architecture Behavioral of Serial_to_Parallel_SR is
    signal S_reg : STD_LOGIC_VECTOR(3 downto 0);
    
begin
    
    process(CLK, RESET)
    begin
        if (RESET = '1') then
            S_reg <= (others => '0');
            
        elsif rising_edge(CLK) then
            
            -- CASE statement for shift control
            case Shift is
                when '1' =>
                    -- Shift enable: shift data right
                    S_reg <= Serial_In & S_reg(3 downto 1);
                
                when others =>
                    -- Hold value when shift disabled
                    S_reg <= S_reg;
            end case;
        end if;
    end process;
    
    S <= S_reg;
    
end Behavioral;
\end{lstlisting}

\subsection{Operation}

\begin{table}[h]
\centering
\begin{tabular}{|c|c|c|c|c|c|c|}
\hline
\textbf{CLK} & \textbf{Shift} & \textbf{Serial In} & \textbf{S[3]} & \textbf{S[2]} & \textbf{S[1]} & \textbf{S[0]} \\
\hline
0 & - & - & 0 & 0 & 0 & 0 \\
\hline
$\uparrow$ 1 & 1 & 1 & 1 & 0 & 0 & 0 \\
\hline
$\uparrow$ 2 & 1 & 0 & 0 & 1 & 0 & 0 \\
\hline
$\uparrow$ 3 & 1 & 1 & 1 & 0 & 1 & 0 \\
\hline
$\uparrow$ 4 & 1 & 1 & 1 & 1 & 0 & 1 \\
\hline
\end{tabular}
\end{table}

\newpage
\section{Comparison and Selection Guide}

\subsection{Comparison Table}

\begin{table}[h]
\centering
\begin{tabular}{|l|c|c|c|}
\hline
\textbf{Feature} & \textbf{IF-ELSE} & \textbf{CASE} & \textbf{WHEN-ELSE} \\
\hline
Context & Sequential & Sequential & Concurrent \\
\hline
Within Process & Yes & Yes & No \\
\hline
Best for & Complex conditions & Selectors & Simple muxing \\
\hline
Nesting & Supports & Limited & Not applicable \\
\hline
Default required & No & Yes & Yes \\
\hline
\end{tabular}
\end{table}

\subsection{When to Use}

\textbf{Use IF-ELSE When:}
\begin{itemize}
    \item Hierarchical or nested conditions
    \item Complex sequential logic
    \item Reset logic in flip-flops/latches
    \item Priority-based logic
\end{itemize}

\textbf{Use CASE When:}
\begin{itemize}
    \item Multiple discrete values
    \item Mutually exclusive branches
    \item Clean, readable code
    \item Multiplexers and decoders
\end{itemize}

\textbf{Use WHEN-ELSE When:}
\begin{itemize}
    \item Combinational logic only
    \item Concurrent assignment
    \item Simple independent conditions
    \item Short and clear code
\end{itemize}

\newpage
\section{Best Practices}

\begin{enumerate}
    \item \textbf{Always cover all cases}
    \begin{itemize}
        \item Add \texttt{when others} to case statements
        \item Provide default value in when-else
    \end{itemize}
    
    \item \textbf{Use meaningful names}
    \begin{itemize}
        \item Clear signal names (Enable, Reset, Serial\_In)
        \item Make intent obvious in code
    \end{itemize}
    
    \item \textbf{Maintain hierarchical structure}
    \begin{itemize}
        \item Keep reset logic separate
        \item Use if-elsif-else for dependent conditions
    \end{itemize}
    
    \item \textbf{Be explicit}
    \begin{itemize}
        \item Don't rely on implicit holds
        \item Make sequential logic explicit
    \end{itemize}
    
    \item \textbf{Choose the right tool}
    \begin{itemize}
        \item Case for selectors (optimal synthesis)
        \item If-else for complex conditions
        \item When-else for simple concurrent logic
    \end{itemize}
\end{enumerate}

\section{Common Mistakes to Avoid}

\begin{enumerate}
    \item \textbf{Missing WHEN OTHERS} - Always provide default case
    \item \textbf{Incomplete sensitivity list} - Include all async signals in process
    \item \textbf{Ignoring priority order} - When-else evaluates top to bottom
    \item \textbf{Invalid concurrency} - Case must be in a process
    \item \textbf{Undefined behavior} - Always assign all signals explicitly
\end{enumerate}

\end{document}
